\documentclass[12pt,a4paper]{article}
\usepackage[T2A]{fontenc}
\usepackage[utf8]{inputenc}
\usepackage[mongolian]{babel}
\usepackage{amsmath}
\usepackage{booktabs}
\usepackage{hyperref}

\title{Монгол хэлний үгийн алдааг автоматаар илрүүлэх нь}
\author{Овог Нэр\\ Их сургууль, Мэдээллийн технологийн тэнхим}
\date{\today}

\begin{document}

\maketitle

\begin{abstract}
Энэхүү өгүүлэлд толь бичигт суурилсан үг шалгагчийн ажиллах зарчим, нарийвчлалыг товч авч үзнэ.
\end{abstract}

\section{Оршил}
Үгийн алдаа нь бичвэрийн чанарт шууд нөлөөлдөг. Иймд алдааг автоматаар илрүүлэх хэрэгсэл чухал~\cite{hunspell}.

\section{Арга зүй}
Үг бүрийг толь бичигтэй харьцуулж, олдоогүй үгэнд хамгийн ойр хувилбарыг санал болгоно. Санал болгох магадлалыг дараах томьёогоор үнэлнэ:
\begin{equation}
  P(w \mid c) = \frac{N(w, c)}{N(c)}
\end{equation}

\section{Үр дүн}
\begin{table}[h]
  \centering
  \begin{tabular}{lr}
    \toprule
    Үзүүлэлт & Утга \\
    \midrule
    Нарийвчлал & 0.94 \\
    Бүрэн байдал & 0.91 \\
    \bottomrule
  \end{tabular}
  \caption{Туршилтын үр дүн}
\end{table}

\section{Дүгнэлт}
Толь бичигт суурилсан арга нь энгийн бөгөөд үр дүнтэй болох нь харагдлаа.

\begin{thebibliography}{9}
\bibitem{hunspell} Hunspell: \url{https://hunspell.github.io}
\end{thebibliography}

\end{document}
